\documentclass[12pt]{article}
\usepackage[utf8]{inputenc}
\usepackage[T1]{fontenc}
\usepackage[margin=1in]{geometry}
\usepackage{setspace}
\usepackage{xcolor}
\usepackage[colorlinks=true, urlcolor=blue!60!black, linkcolor=blue!60!black]{hyperref}
\setlength{\parindent}{0em}
\setlength{\parskip}{0.6em}
\renewcommand{\baselinestretch}{1.0}
\begin{document}
\begin{center}
{\LARGE \textit{Flast v. Cohen}, 392 U.S. 83 (1968)\par}
\vspace{0.5em}
{\large Case Brief --- Two-Page Edition\par}
\end{center}

\section{Memory Jogger}
Taxpayers may challenge congressional spending that violates a specific limit on the spending power, here the Establishment Clause; \textit{Frothingham} narrowed (\href{https://supreme.justia.com/cases/federal/us/392/83/}{Justia}).

\section{Detailed Case Facts}
Florence Flast and six other federal taxpayers sued HEW Secretary Wilbur Cohen to stop grants under the Elementary and Secondary Education Act of 1965 that allegedly financed instruction and textbooks in religious schools, claiming standing solely as taxpayers.

\section{Procedural History}
A three-judge district court dismissed under \textit{Frothingham v.\ Mellon} (1923), over Judge Frankel's dissent; the plaintiffs appealed directly.

\section{Judicial Votes}
Reversed, 8--1. Warren wrote for the Court; Douglas, Stewart, and Fortas concurred; Harlan dissented.

\section{Holding}
Article III does not bar taxpayer suits. A taxpayer has standing by showing (1) a logical link between taxpayer status and an exercise of the Article I, \S\ 8 taxing and spending power, and (2) a nexus to a specific constitutional limit on that power, such as the Establishment Clause.

\section{Analysis of Opinions}
\textbf{Majority:} Standing asks whether the plaintiff has a personal stake assuring ``concrete adverseness''; separation of powers concerns the issues, not the party. Madison's \textit{Memorial and Remonstrance} shows the Establishment Clause limits spending; Mrs.\ Frothingham's Tenth Amendment claim did not. \textbf{Douglas:} overrule \textit{Frothingham}. \textbf{Stewart and Fortas:} confine the holding to the Establishment Clause. \textbf{Harlan:} taxpayers are ``non-Hohfeldian'' plaintiffs, and neither nexus measures a personal stake.

\section{Examples: Future Applications}
\textbf{Same-Side:} a statute appropriating \$50 million for repair grants to churches and synagogues is challengeable. \textbf{Opposite-Side:} an executive gift of surplus property to a seminary (\textit{Valley Forge}) or a religious tax credit (\textit{Winn}) is not. \textbf{Fence-Sitter:} an agency-funded worship component that may be either discretionary spending or part of a statutory grant program.

\section{Critique}
Critics from Harlan to Scalia and Kagan note the arbitrariness of granting standing for direct congressional spending but not for executive spending, land transfers, or tax credits, and \textit{Lujan} repudiated Warren's view that standing does not implicate separation of powers.

\section{Key Quotations}
\begin{quote}
``First, the taxpayer must establish a logical link between that status and the type of legislative enactment attacked. Thus, a taxpayer will be a proper party to allege the unconstitutionality only of exercises of congressional power under the taxing and spending clause of Art. I, \S\ 8, of the Constitution. ... Secondly, the taxpayer must establish a nexus between that status and the precise nature of the constitutional infringement alleged.''
\end{quote}
\noindent --- Chief Justice Warren, Opinion of the Court, 392 U.S. at 102.

\end{document}
